\chapter{Конструкторский раздел}

\section{Требования к разрабатываемому методу}

Разрабатываемый метод построения деревьев составляющих для ограниченного естественного русского языка должен удовлетворять следующим требованиям:
\begin{enumerate}
	\item возможность работы с грамматикой не приведенной к НФХ~---~метод должен успешно обрабатывать правила вида $A \rightarrow B C D$ и $A \rightarrow B$, где в правой части правил располагаются нетерминалы, распространенные в грамматике русского языка, но не удовлетворяющие НФХ;
	\item возможность обработки морфологической омонимии~---~в случае наличия нескольких вариантов морфологического разбора метод должен обработать каждый из них;
	\item возможность учета правил согласования~---~существующие методы работают с грамматикой без учета правил согласования, поскольку данные правила мало распространены в романской и германской языковых группах для которых эти методы разрабатывались; в то же время в русском языке правила согласования важны для грамматически правильного синтаксического анализа~\cite{Arakin_2010,Gak_1983};
	\item возможность определения признаков нетерминальной синтаксической единицы~---~метод должен определять морфологические признаки синтаксической единицы на основе морфологических признаков единиц-потомков;
	\item  возможность обработки пунктуации~---~метод должен обрабатывать знаки пунктуации внутри предложения, в связи с их ролью в обособлении синтаксических единиц;
	\item возможность получения всех возможных деревьев составляющих~---~метод должен предоставлять все возможные варианты выделения синтаксических единиц в рамках предложения.
\end{enumerate}

Для удовлетворения перечисленным требованиям, разработаны следующие решения:
\begin{enumerate}
	\item приведение правил грамматики к правилам удовлетворяющим НФХ в ходе работы метода, путем ввода дополнительных вспомогательных нетерминалов, которые в дальнейшем удаляются из результирующего дерева составляющих; потомки таких нетерминальных вершин, становятся потомками предка нетерминала;
	\item внесение на начальном этапе формирования таблицы составляющих по CYK всех возможных наборов морфологических характеристик каждой словоформы;
	\item объединение синтаксических единиц в ходе заполнения таблицы составляющих не только по правилам грамматики, но и правилам согласования, перечисленным в приложении В;
	\item определение морфологических характеристик не листовой вершины в дереве составляющих по морфологическим характеристикам ее потомков на основе вида синтаксической единицы, которой соответствует вершина, и правил согласования, перечисленных в приложении B;
	\item обработка знаков пунктуации как отдельных терминалов;
	\item рекурсивное восстановление дерева составляющих по полученной таблице составляющих с сохранением всех вариантов.
\end{enumerate}

Разрабатываемый метод имеет следующие ограничения:
\begin{itemize}
	\item работа с предложениями на ограниченном естественном русском языке, который задается наборами правил грамматики и согласования, описанными в приложениях Б и В соответственно;
	\item предложения не удовлетворяющие описанным правилам грамматики или согласования не обрабатываются;
	\item невозможность полноценной обработки непроективных конструкций;
	\item к обрабатываемым знакам пунктуации относятся запятые, двоеточия, точки с запятой; прочие знаки пунктуации не обрабатываются методом.
\end{itemize}

\section{Используемые структуры и форматы данных}

Для обработки данных разрабатываемым методом, необходимо выбрать используемые им структуры и форматы представления данных. В данном разделе описываются представления данных с которыми работает метод на различных этапах и используемые для этого структуры данных.

Для представления входной токенизированной последовательности перед обработкой морфологическим анализатором используется массив строк, содержащий словоформы входного предложения и обрабатываемые знаки препинания. Такое представление позволяет за линейное время обработать все словоформы предложения.

Последовательность словоформ получивших морфологические характеристики, а так же обрабатываемые знаки препинания, представляется в виде массива ассоциативных массивов, в которых ключами являются названия морфологических характеристик, значениями~---~значения морфологических характеристик. Такой подход позволяет обработать все возможные морфологические разборы за линейное время, а получение в конкретном разборе определенной характеристики осуществляется за $O(1)$ при реализации ассоциативного массива с помощью хеш-таблицы~\cite{cormen13}.

Таблица составляющих с которой работает алгоритм CYK представлена двумерным массивом, каждая ячейка которого~---~ассоциативный массив, в котором ключ~---~название синтаксической единицы, значение~---~множество возможных наборов морфологических характеристик данной единицы. Такой подход позволяет для каждого подотрезка входного предложения сохранить более одной возможной синтаксической единицы, а для каждой такой единицы в свою очередь более одного варианта набора морфологических характеристик. Множественные варианты необходимы вследствие морфологической и синтаксической омонимии, которые приводят к различным вариантам разборов подотрезков данного отрезка и, в простейшем случае, к различным вариантам разбора одной словоформы. Множество наборов морфологических характеристик в свою очередь состоит из неизменяемых множеств пар строковых значений, определяющих название и значение морфологических характеристик. Данное представление позволяет для каждой синтаксической единицы поддерживать изменяемое множество уникальных неизменяемых наборов морфологических характеристик во избежание повторяющихся наборов таких характеристик и проверки наличия определенного набора среди уже имеющихся за $O(1)$. При этом каждый набор представлен неизменяемым множеством, чтобы не допустить дублирование характеристик в одном наборе.

Грамматика представляется в виде ассоциативного массива, ключем в котором является название синтаксической единицы (нетерминал), значением~---~массив возможных разборов данной единицы. Каждый возможный разбор представляется массивом названий синтаксических единиц-потомков нетерминала-ключа. Такое представление позволяет за $O(1)$ определить потомков нетерминала. Однако, поскольку алгоритм CYK производит разбор снизу-вверх~\cite{jurafsky2026speech}, данное представление следует инвертировать, для получение нетерминала, который может быть составлен из двух данных синтаксических единиц. Таким образом, в инвертированном представлении грамматика является ассоциативным массивом, в котором ключами являются пары синтаксических единиц, а значениями~---~массив всех выводимых из них нетерминалов.

Для представления получаемого в результате работы метода дерева используется формат используется ассоциативный массив содержащий в качестве значений:
\begin{itemize}
	\item имя синтаксической единицы соответствующей данной вершине;
	\item ассоциативный массив такой же структуры, описывающий потомков данной вершины;
	\item словоформу, соответствующую данной вершине, если она является листовой.
\end{itemize} 

Для хранения грамматики грамматических правил и получаемых деревьев синтаксического разбора могут быть использованы файлы формата JSON, так как они соответствуют структуре ассоциативного массива, могут быть сереализованы в такие массивы и десерализованы из них средствами большинства языков программирования~\cite{rfc8259}.

\section{Описание ключевых этапов метода}